%% Question bank chapter file - generated by bank_tools.py build (do not edit while the bank is still being built:
%% it is regenerated from the per-test data in Question Bank\data). Every question is one bankq block:
%%   \begin{bankq}[id=..., source={...}, paper=.., qno=.., type=..., subtopic={...}, answer={...}, answer_src=key|solved,
%%                 working={...}, difficulty=1-10, diff_parts={K S I C F}, diff_note={...},
%%                 group=<passage id or empty>, also_in={duplicate copies in other tests}]
%% difficulty = min(10, K+S+I+C+F-1): K concepts 1-3, S steps 1-4, I insight 0-2, C computation 0-2, F format 0-1
%%   ...question LaTeX, options, figures (\bfig / \bfigi -> img/)...
%%   \end{bankq}
%% Passages shared by several questions are bankpassage blocks; their questions carry group=<passage id>.
\documentclass[10pt]{article}
\usepackage[a4paper,margin=1.8cm]{geometry}
\usepackage[T1]{fontenc}
\usepackage{amsmath,amssymb,mathtools}
\usepackage{graphicx}
\usepackage[export]{adjustbox}
\usepackage{xcolor,needspace,keyval}
\usepackage{enumitem}
\usepackage{booktabs,array,tabularx}
\usepackage[hidelinks]{hyperref}
\setlength{\parindent}{0pt}
\setlength{\parskip}{3pt}
\makeatletter
\define@key{bq}{id}{\def\bq@id{#1}}\define@key{bq}{source}{\def\bq@src{#1}}\define@key{bq}{paper}{\def\bq@paper{#1}}
\define@key{bq}{qno}{\def\bq@qno{#1}}\define@key{bq}{type}{\def\bq@type{#1}}\define@key{bq}{subtopic}{\def\bq@sub{#1}}
\define@key{bq}{answer}{\def\bq@ans{#1}}\define@key{bq}{answer_src}{\def\bq@asrc{#1}}\define@key{bq}{working}{\def\bq@work{#1}}
\define@key{bq}{group}{\def\bq@group{#1}}\define@key{bq}{also_in}{\def\bq@also{#1}}
\define@key{bq}{difficulty}{\def\bq@diff{#1}}\define@key{bq}{diff_parts}{\def\bq@dparts{#1}}\define@key{bq}{diff_note}{\def\bq@dnote{#1}}
\newenvironment{bankq}[1][]{\def\bq@work{}\def\bq@also{}\def\bq@group{}\def\bq@diff{?}\def\bq@dparts{}\def\bq@dnote{}\setkeys{bq}{#1}%
  \par\needspace{4\baselineskip}\medskip\noindent{\small\textbf{\bq@id}\hfill\color{gray}[\bq@src, Paper \bq@paper, Q\bq@qno]
  \bq@type\ $\cdot$ \bq@sub\ $\cdot$ \textbf{D\bq@diff}}\par\nopagebreak\relax}
  {\par{\small\color{teal}Answer: \textbf{\bq@ans}\ifx\bq@work\@empty\else\ (solved: \bq@work)\fi
  \ifx\bq@also\@empty\else\par\color{gray}Also in: \bq@also\fi}\par}
\define@key{bp}{id}{\def\bp@id{#1}}\define@key{bp}{source}{\def\bp@src{#1}}\define@key{bp}{questions}{\def\bp@qs{#1}}
\newenvironment{bankpassage}[1][]{\setkeys{bp}{#1}\par\needspace{4\baselineskip}\bigskip\noindent
  {\small\textbf{Passage \bp@id}\hfill\color{gray}[\bp@src; questions \bp@qs]}\par\nopagebreak\relax\itshape}{\par}
\makeatother
\newcommand{\bfig}[2]{\par\smallskip{\centering\includegraphics[width=#1pt,max width=0.62\linewidth,max height=0.32\textheight]{img/#2}\par}\smallskip}
\newcommand{\bfigi}[2]{\includegraphics[width=#1pt,max width=\linewidth,max height=0.32\textheight,valign=c]{img/#2}}
\newcommand{\bmissing}[1]{\textit{[figure missing in the source conversion: \detokenize{#1}]}}
\newenvironment{qoptions}{\begin{itemize}[leftmargin=1.8em,itemsep=1pt,topsep=2pt,label={}]}{\end{itemize}}
\newcommand{\opt}[2]{\item[(#1)] #2}
\providecommand{\fwstart}{}\providecommand{\fwend}{}\providecommand{\qfullsplit}{}
\newcolumntype{L}[1]{>{\raggedright\arraybackslash}p{#1}}
\newcolumntype{C}[1]{>{\centering\arraybackslash}p{#1}}
\newcolumntype{R}[1]{>{\raggedleft\arraybackslash}p{#1}}
\providecommand{\sectionhdr}[1]{}\providecommand{\subsectionhdr}[1]{}


\begin{document}
{\Large\bfseries Periodic Table --- Question Bank}\par\smallskip
{\small Chemistry chapter C29. Questions from all converted test papers, grouped by source series. Answers: printed key where it matches the paper, otherwise solved (answer\_src=solved, with working). Exact repeats across papers appear once, with also\_in listing the other copies.}\par\medskip\hrule

\section*{Score-II papers}
\begin{bankq}[id=SII02-P2-CHEM-S1-Q02, source={Score-II - Test 02 (28 Mar 2021)}, paper=2, qno=2, type=multi, subtopic={Periodic properties}, answer={B,D}, answer_src=key, difficulty=8, diff_parts={K3 S1 I2 C2 F1}, diff_note={concepts: standard problem in this chapter, periodic properties, apply the relevant results | steps: apply the relevant results; compute}]
IE$_1$ (kJ/mole) of $2^{\text{nd}}$ period elements are mentioned below with some blank/vacant data spaces:

\begin{center}
\renewcommand{\arraystretch}{1.2}
\begin{tabular}{cccccccc}
Li & Be & B & C & N & O & F & Ne\\
520 & ``P'' & 800 & ``Q'' & 1402 & 1314 & ``R'' & 2080\\
\end{tabular}
\renewcommand{\arraystretch}{1.15}
\end{center}

Select CORRECT number of statement(s) regarding above data table. (Consider all processes are gas phased).

\begin{qoptions}
\opt{A}{$|$EA$|$ of N$^+$ is more than 1402 kJ/mol}
\opt{B}{Value of ``R'' is more than 1402 kJ/mol}
\opt{C}{Value of ``P'' is more than $|$EA$|$ of C$^+$}
\opt{D}{Q $>$ 800 kJ/mol}
\end{qoptions}
\end{bankq}
\begin{bankq}[id=SII07-P1-CHEM-S1-Q11, source={Score-II - Test 07 (10 Jul 2022)}, paper=1, qno=11, type=single, subtopic={Trends}, answer={D}, answer_src=key, difficulty=4, diff_parts={K2 S1 I1 C1 F0}, diff_note={concepts: standard problem in this chapter, trends | steps: apply the relevant results; compute}]
Select the incorrect order of comparisons for the indicated properties :
\begin{qoptions}
\opt{A}{$\mathrm{Cu^{+2}(aq.)>Cu^{+}(aq.)}$ : stability}
\opt{B}{$\mathrm{Cr_2O_7^{2-}<MnO_4^{-}}$ : oxidising power}
\opt{C}{$\mathrm{Cr^{+2}>Fe^{+2}}$ : reducing power}
\opt{D}{C $>$ Fe $>$ Lanthanoid metal : percentage in mischmetal}
\end{qoptions}
\end{bankq}
\begin{bankq}[id=SII09-P1-CHEM-S1-Q03, source={Score-II - Test 09 (12 Feb 2023)}, paper=1, qno=3, type=multi, subtopic={Trends}, answer={A,B,D}, answer_src=key, difficulty=8, diff_parts={K3 S1 I2 C2 F1}, diff_note={concepts: standard problem in this chapter, trends, apply the relevant results | steps: apply the relevant results; compute}]
Which of the following is/are correct order of the mentioned property
\begin{qoptions}
\opt{A}{Acidic Strength:\\ \bfigi{164.9}{SII09_p1_chem_s1_q03_1.png}}
\opt{B}{Dipole moment:\\ \bfigi{129.6}{SII09_p1_chem_s1_q03_2.png}}
\opt{C}{Heat of hydrogenation:\\ \bfigi{76.8}{SII09_p1_chem_s1_q03_3.png}}
\opt{D}{Basic strength:\\ \bfigi{176.1}{SII09_p1_chem_s1_q03_4.png}}
\end{qoptions}
\end{bankq}
\begin{bankq}[id=SII09-P1-CHEM-S1-Q04, source={Score-II - Test 09 (12 Feb 2023)}, paper=1, qno=4, type=multi, subtopic={Periodic properties}, answer={B,D}, answer_src=key, difficulty=8, diff_parts={K3 S1 I2 C2 F1}, diff_note={concepts: standard problem in this chapter, periodic properties, apply the relevant results | steps: apply the relevant results; compute}]
The INCORRECT fact(s) regarding the modern periodic table is/are:
\begin{qoptions}
\opt{A}{The element that has the maximum number of electrons in the last shell belongs to the d-block}
\opt{B}{The atom that has the highest electron affinity is fluorine.}
\opt{C}{Element having electronic configuration \resizebox{\linewidth}{!}{$1s^22s^22p^63s^23p^63d^{10}4s^24p^64d^{10}4f^{14}5s^25p^65d^{10}6s^26p^66d^27s^2$} belongs to the f-block}
\opt{D}{The ionisation energy increases continuously on moving left to right in all the periods.}
\end{qoptions}
\end{bankq}
\begin{bankq}[id=SII13-P1-CHEM-S1-Q11, source={Score-II - Test 13 (06 Feb 2024)}, paper=1, qno=11, type=single, subtopic={Periodic properties}, answer={B}, answer_src=key, difficulty=4, diff_parts={K2 S1 I1 C1 F0}, diff_note={concepts: standard problem in this chapter, periodic properties | steps: apply the relevant results; compute}]
The ionization potential of nitrogen is more than that of oxygen because of :-
\begin{qoptions}
\opt{A}{Greater attraction of electrons by the nucleus}
\opt{B}{Extra stability of the half filled p-orbitals}
\opt{C}{Smaller size of nitrogen}
\opt{D}{More penetration effect}
\end{qoptions}
\end{bankq}
\begin{bankq}[id=SII13-P1-CHEM-S1-Q12, source={Score-II - Test 13 (06 Feb 2024)}, paper=1, qno=12, type=single, subtopic={Periodic properties}, answer={D}, answer_src=key, difficulty=4, diff_parts={K2 S1 I1 C1 F0}, diff_note={concepts: standard problem in this chapter, periodic properties | steps: apply the relevant results; compute}]
The option with only amphoteric oxides is :-
\begin{qoptions}
\opt{A}{$Cr_2O_3$, NO, SnO, PbO}
\opt{B}{NO, $B_2O_3$, PbO, $SnO_2$}
\opt{C}{$Cr_2O_3$, BaO, SnO, $SnO_2$}
\opt{D}{ZnO, $Al_2O_3$, PbO, $PbO_2$}
\end{qoptions}
\end{bankq}
\begin{bankq}[id=SII15-P1-CHEM-S1-Q15, source={Score-II - Test 15 (09 Feb 2024)}, paper=1, qno=15, type=single, subtopic={Trends}, answer={B}, answer_src=key, difficulty=4, diff_parts={K2 S1 I1 C1 F0}, diff_note={concepts: standard problem in this chapter, trends | steps: apply the relevant results; compute}]
Which of the following is correct order of solubility of given salts in aq solution?
\begin{qoptions}
\opt{A}{Be(NO$_3$)$_2$ $<$ Mg(NO$_3$)$_2$ $<$ Ca(NO$_3$)$_2$ $<$ Sr(NO$_3$)$_2$ $<$ Ba(NO$_3$)$_2$}
\opt{B}{Li$_2$CO$_3$ $<$ Na$_2$CO$_3$ $<$ K$_2$CO$_3$ $<$ Rb$_2$CO$_3$ $<$ Cs$_2$CO$_3$}
\opt{C}{LiF $>$ NaF $>$ KF $>$ RbF $>$ CsF}
\opt{D}{AgF $<$ AgCl $<$ AgBr $<$ AgI}
\end{qoptions}
\end{bankq}
\begin{bankq}[id=SII15-P2-CHEM-S2-Q08, source={Score-II - Test 15 (09 Feb 2024)}, paper=2, qno=8, type=multi, subtopic={Trends}, answer={A,C,D}, answer_src=key, difficulty=8, diff_parts={K3 S1 I2 C2 F1}, diff_note={concepts: standard problem in this chapter, trends, apply the relevant results | steps: apply the relevant results; compute}]
Which of the following order is/are correct?
\begin{qoptions}
\opt{A}{Bond energy : Cl$_2$ $>$ Br$_2$ $>$ F$_2$ $>$ I$_2$}
\opt{B}{Bond strength : $-\underset{|}{\mathrm N}-\underset{|}{\mathrm N}- \;>\; -\underset{|}{\mathrm P}-\underset{|}{\mathrm P}-$}
\opt{C}{Number of $\pi$-bonds : HOCl $<$ HClO$_2$ $<$ HClO$_3$ $<$ HClO$_4$}
\opt{D}{Bond angle : NO$_2^+$ $>$ NO$_2$ $>$ NO$_2^-$}
\end{qoptions}
\end{bankq}
\begin{bankq}[id=SII17-P1-CHEM-S2-Q07, source={Score-II - Test 17 (13 Feb 2024)}, paper=1, qno=7, type=numerical, subtopic={Periodic properties}, answer={4}, answer_src=solved, working={Correct: (b) Zn<Cd<Hg, (c) B>Tl>Ga>Al>In, (d) S>Se>Te>O, (e) Ni<Cu<Zn. Wrong: a, f, g. Total 4.}, difficulty=8, diff_parts={K3 S1 I2 C2 F1}, diff_note={concepts: standard problem in this chapter, periodic properties, apply the relevant results | steps: apply the relevant results; compute}, also_in={Score-II - Test 19 (20 Feb 2024) P1 Q3}]
How many of following orders are correct?\\
(a) Order of IE : O$^{-2}$ $>$ O$^+$ $>$ O $>$ O$^-$\\
(b) Order of EN : Zn $<$ Cd $<$ Hg\\
(c) Order of IE : B $>$ Tl $>$ Ga $>$ Al $>$ In\\
(d) Order of EA : S $>$ Se $>$ Te $>$ O\\
(e) Order of atomic radius : Ni $<$ Cu $<$ Zn\\
(f) Order of ionic radius : Na$^+$ $>$ Mg$^{+2}$ $>$ Li$^+$ $>$ Be$^{+2}$\\
(g) Order of IE$_2$ : N $<$ O $<$ F $<$ Ne
\end{bankq}
\begin{bankq}[id=SII19-P1-CHEM-S1-Q03, source={Score-II - Test 19 (20 Feb 2024)}, paper=1, qno=3, type=single, subtopic={Periodic properties}, answer={C}, answer_src=key, difficulty=4, diff_parts={K2 S1 I1 C1 F0}, diff_note={concepts: standard problem in this chapter, periodic properties | steps: apply the relevant results; compute}]
Which element is having lowest melting and boiling point?
\begin{qoptions}
\opt{A}{Ti}
\opt{B}{Cu}
\opt{C}{Zn}
\opt{D}{Mn}
\end{qoptions}
\end{bankq}
\begin{bankq}[id=SII22-P1-CHEM-S1-Q07, source={Score-II - Test 22 (14 Mar 2024)}, paper=1, qno=7, type=single, subtopic={Periodic properties}, answer={A}, answer_src=key, difficulty=4, diff_parts={K2 S1 I1 C1 F0}, diff_note={concepts: standard problem in this chapter, periodic properties | steps: apply the relevant results; compute}]
To which of the following atom, the attachment of electron is most difficult ?
\begin{qoptions}
\opt{A}{Radon}
\opt{B}{Nitrogen}
\opt{C}{Oxygen}
\opt{D}{Radium}
\end{qoptions}
\end{bankq}
\begin{bankq}[id=SII22-P1-CHEM-S2-Q05, source={Score-II - Test 22 (14 Mar 2024)}, paper=1, qno=5, type=numerical, subtopic={Periodic properties}, answer={3}, answer_src=key, difficulty=8, diff_parts={K3 S1 I2 C2 F1}, diff_note={concepts: standard problem in this chapter, periodic properties, apply the relevant results | steps: apply the relevant results; compute}]
Count basic hydroxides NaOH, ClOH, KOH, $Ca(OH)_2$, $Zn(OH)_2$, $(OH)_2SO_2$, $HO(NO_2)$, $Al(OH)_3$
\end{bankq}
\begin{bankq}[id=SII24-P1-CHEM-S1-Q02, source={Score-II - Test 24 (17 Mar 2024)}, paper=1, qno=2, type=single, subtopic={Trends}, answer={B}, answer_src=key, difficulty=4, diff_parts={K2 S1 I1 C1 F0}, diff_note={concepts: standard problem in this chapter, trends | steps: apply the relevant results; compute}]
Which of the following order is correct?
\begin{qoptions}
\opt{A}{B $<$ Al $<$ Ga $<$ In $<$ T$\ell$ (Atomic radius)}
\opt{B}{B$^{+3}$ $<$ Al$^{+3}$ $<$ Ga$^{+3}$ $<$ In$^{+3}$ $<$ Tl$^{+3}$ (Ionic radius)}
\opt{C}{B $>$ Ga $>$ Al $>$ In $>$ T$\ell$ (Ionization energy)}
\opt{D}{B $>$ Al $>$ Ga $>$ In $>$ T$\ell$ (Electronegativity)}
\end{qoptions}
\end{bankq}
\begin{bankq}[id=SII26-P1-CHEM-S1-Q03, source={Score-II - Test 26 (20 Mar 2024)}, paper=1, qno=3, type=single, subtopic={Periodic properties}, answer={D}, answer_src=key, difficulty=4, diff_parts={K2 S1 I1 C1 F0}, diff_note={concepts: standard problem in this chapter, periodic properties | steps: apply the relevant results; compute}]
Choose the correct statements
\begin{qoptions}
\opt{A}{Cl$_2$ is strongest oxidising halogen since Cl$_2$ has highest EA}
\opt{B}{BDE order : F$_2$ $>$ Cl$_2$ $>$ Br$_2$ $>$ I$_2$}
\opt{C}{Maximum number of atoms in a plane for I$_2$Cl$_6$ and B$_2$H$_6$ are same.}
\opt{D}{BF$_3$ and BrF$_3$ both are planar.}
\end{qoptions}
\end{bankq}
\begin{bankq}[id=SII27-P1-CHEM-S1-Q20, source={Score-II - Test 27 (23 Mar 2024)}, paper=1, qno=20, type=single, subtopic={Trends}, answer={C}, answer_src=key, difficulty=4, diff_parts={K2 S1 I1 C1 F0}, diff_note={concepts: standard problem in this chapter, trends | steps: apply the relevant results; compute}]
Select the incorrect statement.
\begin{qoptions}
\opt{A}{There is gradual change from basic V$_2$O$_3$ to less basic V$_2$O$_4$ and to amphoteric V$_2$O$_5$}
\opt{B}{V$_2$O$_4$ dissolves in acids to give VO$^{2+}$ salts}
\opt{C}{V$_2$O$_5$ reacts with alkalies as well as acids to give V O$_4^{+}$ and V O$_4^{3-}$ respectively}
\opt{D}{V$_2$O$_5$ is amphoteric though mainly acidic}
\end{qoptions}
\end{bankq}
\begin{bankq}[id=SII29-P1-CHEM-S2-Q07, source={Score-II - Test 29 (26 Mar 2024)}, paper=1, qno=7, type=multi, subtopic={Trends}, answer={A,C,D}, answer_src=key, difficulty=8, diff_parts={K3 S1 I2 C2 F1}, diff_note={concepts: standard problem in this chapter, trends, apply the relevant results | steps: apply the relevant results; compute}]
Which of the following statement is correct?

\begin{qoptions}
\opt{A}{The melting point of hydrides of Group 15 : NH$_3$ > SbH$_3$ > AsH$_3$ > PH$_3$}
\opt{B}{$\Delta H_{eg}$ kJ mol$^{-1}$ : S > Se > O > Te > Po}
\opt{C}{Melting point : B > Al > Tl > In > Ga}
\opt{D}{Bond enthalpy : C--C > Si--Si > Ge--Ge > Sn--Sn}
\end{qoptions}
\end{bankq}
\begin{bankq}[id=SII30-P1-CHEM-S1-Q12, source={Score-II - Test 30 (27 Mar 2024)}, paper=1, qno=12, type=single, subtopic={Trends}, answer={D}, answer_src=key, difficulty=4, diff_parts={K2 S1 I1 C1 F0}, diff_note={concepts: standard problem in this chapter, trends | steps: apply the relevant results; compute}]
In case of isoelectronic species the size of F$^{-}$, Ne and Na$^{+}$ is affected by:
\begin{qoptions}
\opt{A}{Principal quantum number ($n$)}
\opt{B}{None of the factors because their size is the same}
\opt{C}{Electron-electron interaction in the outer orbitals}
\opt{D}{Nuclear charge ($z$)}
\end{qoptions}
\end{bankq}
\begin{bankq}[id=SII31-P1-CHEM-S1-Q02, source={Score-II - Test 31 (29 Mar 2024)}, paper=1, qno=2, type=single, subtopic={Trends}, answer={B}, answer_src=key, difficulty=4, diff_parts={K2 S1 I1 C1 F0}, diff_note={concepts: standard problem in this chapter, trends | steps: apply the relevant results; compute}]
Select the \textbf{CORRECT} order according to give property.
\begin{qoptions}
\opt{A}{Lewis acidic nature : BCl$_3$ $<$ AlCl$_3$ $<$ GaCl$_3$}
\opt{B}{Response for magnet : CaSO$_4$ $<$ FeSO$_4$ $<$ MnSO$_4$}
\opt{C}{Reducing property : AsH$_3$ $<$ PH$_3$ $<$ NH$_3$}
\opt{D}{Solubility : AgF $<$ AgCl $<$ AgBr}
\end{qoptions}
\end{bankq}
\begin{bankq}[id=SII31-P1-CHEM-S2-Q02, source={Score-II - Test 31 (29 Mar 2024)}, paper=1, qno=2, type=single, subtopic={Trends}, answer={2}, answer_src=key, difficulty=7, diff_parts={K3 S1 I2 C2 F0}, diff_note={concepts: standard problem in this chapter, trends, apply the relevant results | steps: apply the relevant results; compute}]
Consider the following order according to indicate property.\\
(i) B$-$F bond length : BF$_3$ $<$ BF$_2$OH\\
(ii) Dipole moment : CHCl$_3$ $<$ CH$_3$Cl\\
(iii) Acidic strength : SO$_3$ $<$ SeO$_3$\\
(iv) Number of d-orbital , participate in hybridisation : BrF$_5$ $<$ PF$_5$\\
Find the total number of correct orders?
\end{bankq}
\begin{bankq}[id=SII35-P1-CHEM-S2-Q07, source={Score-II - Test 35 (04 Feb 2025)}, paper=1, qno=7, type=multi, subtopic={Periodic properties}, answer={B,C,D}, answer_src=key, difficulty=8, diff_parts={K3 S1 I2 C2 F1}, diff_note={concepts: standard problem in this chapter, periodic properties, apply the relevant results | steps: apply the relevant results; compute}]
The possible type of ions that can be produced by oxides of Vanadium on treatment with acidic or alkaline solutions is/are.
\begin{qoptions}
\opt{A}{VO$^+$}
\opt{B}{VO$^{2+}$}
\opt{C}{VO$_2^+$}
\opt{D}{VO$_4^{3-}$}
\end{qoptions}
\end{bankq}
\begin{bankq}[id=SII38-P1-CHEM-S1-Q03, source={Score-II - Test 38 (12 Feb 2025)}, paper=1, qno=3, type=single, subtopic={Trends}, answer={C}, answer_src=key, difficulty=4, diff_parts={K2 S1 I1 C1 F0}, diff_note={concepts: standard problem in this chapter, trends | steps: apply the relevant results; compute}]
Which of the following statement/relationship is incorrect?
\begin{qoptions}
\opt{A}{Reducing power in aq. medium : Li $>$ K $>$ Na}
\opt{B}{Solubility in water : Ca(OH)$_2$ $<$ Sr(OH)$_2$ $<$ Ba(OH)$_2$}
\opt{C}{NO is obtained as one of the product when Sn reacts with conc. HNO$_3$}
\opt{D}{Lithium carbide produces acetylene gas on hydrolysis whereas aluminium carbide produces methane gas on hydrolysis}
\end{qoptions}
\end{bankq}
\begin{bankq}[id=SII43-P1-CHEM-S1-Q14, source={Score-II - Test 43 (23 Mar 2025)}, paper=1, qno=14, type=single, subtopic={Trends}, answer={D}, answer_src=key, difficulty=4, diff_parts={K2 S1 I1 C1 F0}, diff_note={concepts: standard problem in this chapter, trends | steps: apply the relevant results; compute}]
Consider Be, B, C, N, O, F, Cl, S atoms and identify atom which has maximum E.A. and atom which has lowest ionization energy.\\
What is the shape of binary molecule formed by them.
\begin{qoptions}
\opt{A}{Linear}
\opt{B}{Tetrahedral}
\opt{C}{T-pyramidal}
\opt{D}{Triangular planar}
\end{qoptions}
\end{bankq}

\section*{Score Advanced papers}
\begin{bankq}[id=SA04-P1-CHEM-S1-Q04, source={Score Advanced - Test 04 (04 Jul 2021)}, paper=1, qno=4, type=multi, subtopic={Periodic properties}, answer={A,B}, answer_src=solved, working={A and B are correct; C and D are reversed}, difficulty=6, diff_parts={K2 S2 I1 C1 F1}, diff_note={concepts: periodic order of size and electron affinity and ionisation energy, periodic properties | steps: Al 143 pm exceeds Ga 135 pm; P3- is much larger than Ca2+; Ne has no positive electron affinity; B exceeds Li in ionisation energy}]
Which of the following order(s) is/are correct?
\begin{qoptions}
\opt{A}{Al $>$ Ga (Atomic Size)}
\opt{B}{P$^{-3}$ $>$ Ca$^{+2}$ (Ionic Size)}
\opt{C}{Ne $>$ Al (Electron Affinity)}
\opt{D}{Li $>$ B (First Ionisation Energy)}
\end{qoptions}
\end{bankq}
\begin{bankq}[id=SA04-P2-CHEM-S1-Q09, source={Score Advanced - Test 04 (04 Jul 2021)}, paper=2, qno=9, type=single, subtopic={Periodic properties}, answer={C}, answer_src=solved, working={Al (1.5) < Si (1.8), so C is wrong}, difficulty=5, diff_parts={K2 S2 I1 C1 F0}, diff_note={concepts: order of ionisation energy and size and electronegativity, periodic properties | steps: Sc>Y>La; lanthanoid contraction; Al is less electronegative than Si}]
Which of the following order is NOT correct, against mentioned properties?
\begin{qoptions}
\opt{A}{Sc $>$ Y $>$ La (Ionisation energy)}
\opt{B}{Hf$^{2+}$ $<$ Hf $\simeq$ Zr (size)}
\opt{C}{A$\ell$ $>$ Si (Electronegativity)}
\opt{D}{Ga $>$ A$\ell$ (Ionisation energy)}
\end{qoptions}
\end{bankq}
\begin{bankq}[id=SA04-P2-CHEM-S1-Q13, source={Score Advanced - Test 04 (04 Jul 2021)}, paper=2, qno=13, type=single, subtopic={Periodic properties}, answer={C}, answer_src=solved, working={O- + e -> O2- absorbs energy}, difficulty=5, diff_parts={K2 S1 I2 C1 F0}, diff_note={concepts: energy absorbed in electron processes, periodic properties | steps: adding a second electron to O- is endothermic; the other three release energy}]
In which of the following process energy is absorbed?
\begin{qoptions}
\opt{A}{S$^{2-}_{(g)}$ $\rightarrow$ S$^{-}_{(g)}$ + e$^-$}
\opt{B}{N$^{-}_{(g)}$ $\rightarrow$ N$_{(g)}$ + e$^-$}
\opt{C}{O$^{-}_{(g)}$ + e$^-$ $\rightarrow$ O$^{2-}_{(g)}$}
\opt{D}{Na$^{+}_{(g)}$ + e$^-$ $\rightarrow$ Na$_{(g)}$}
\end{qoptions}
\end{bankq}
\begin{bankq}[id=SA04-P2-CHEM-S2-Q06, source={Score Advanced - Test 04 (04 Jul 2021)}, paper=2, qno=6, type=numerical, subtopic={Periodic properties}, answer={19}, answer_src=solved, working={12+7=19}, difficulty=4, diff_parts={K1 S1 I1 C1 F1}, diff_note={concepts: group and period of element 112 | steps: Uub is copernicium, group 12, period 7}]
Give the sum of group number and period of element ``Uub''.
\end{bankq}
\begin{bankq}[id=SA05-P1-CHEM-S1-Q01, source={Score Advanced - Test 05 (15 Jul 2021)}, paper=1, qno=1, type=single, subtopic={Periodic properties}, answer={B}, answer_src=key, difficulty=5, diff_parts={K2 S2 I1 C1 F0}, diff_note={concepts: order not agreeing with the property, periodic properties | steps: ionic sizes; second ionisation enthalpies; electron gain; metallic radius}]
In which of the following options order of arrangement does not agree with the variation in property indicated against it.
\begin{qoptions}
\opt{A}{Al$^{+3}$ $<$ Mg$^{+2}$ $<$ Na$^+$ $<$ F$^-$ (Increasing ionic size)}
\opt{B}{B$^+$ $<$ C$^+$ $<$ N$^+$ $<$ O$^+$ (Ionisation enthalpy)}
\opt{C}{I $<$ Br $<$ F $<$ Cl (Increasing $|\Delta$Heg$|$)}
\opt{D}{Li $<$ Na $<$ K $<$ Rb (Increasing metallic radius)}
\end{qoptions}
\end{bankq}
\begin{bankq}[id=SA06-P2-CHEM-S1-Q02, source={Score Advanced - Test 06 (04 Sep 2021)}, paper=2, qno=2, type=multi, subtopic={Periodic properties}, answer={B,C}, answer_src=key, difficulty=9, diff_parts={K3 S2 I2 C2 F1}, diff_note={concepts: compounds possible from ionisation energies, periodic properties, large jump after the third ionisation puts W in group 13 | steps: large jump after the third ionisation puts W in group 13; jump after the second puts X in group 2; Y is group 16 and Z is group 17}]
Consider the following values of I.E.(eV) for elements W and X:\par
{\footnotesize\setlength{\tabcolsep}{3pt}\noindent\begin{tabular}{llll l}
Element & I.E.$_1$ & I.E.$_2$ & I.E.$_3$ & I.E.$_4$ \\
W & 10.5 & 15.5 & 24.9 & 79.8 \\
X & 8 & 14.8 & 78.9 & 105.8 \\
\end{tabular}}\par
Other two element Y and Z have outer electronic configuration ns$^2$np$^4$ and ns$^2$np$^5$ respectively. Then according to given information which of the following compound(s) is/are not possible.
\begin{qoptions}
\opt{A}{W$_2$Y$_3$}
\opt{B}{X$_2$Y$_3$}
\opt{C}{WZ$_2$}
\opt{D}{XZ$_2$}
\end{qoptions}
\end{bankq}
\begin{bankq}[id=SA08-P1-CHEM-S2-Q08, source={Score Advanced - Test 08 (22 Apr 2023)}, paper=1, qno=8, type=multi, subtopic={Periodic properties}, answer={A,B,C,D}, answer_src=key, difficulty=6, diff_parts={K2 S2 I1 C1 F1}, diff_note={concepts: incorrect periodic statements, periodic properties | steps: second ionisation energy order; electron affinity of N is lowest; hafnium is in d-block}]
Which of the following statements is/are incorrect?
\begin{qoptions}
\opt{A}{Correct order of 2$^{\text{nd}}$ I.E. is O $>$ N $>$ C $>$ B}
\opt{B}{Correct order of EA is O $>$ C $>$ B $>$ Li $>$ N}
\opt{C}{Element with Z = 72 belongs to 6$^{\text{th}}$ period \& p-block}
\opt{D}{I.E. order for 13$^{\text{th}}$ group element is B $>$ T$\ell$ $>$ Al $>$ Ga $>$ In}
\end{qoptions}
\end{bankq}
\begin{bankq}[id=SA17-P1-CHEM-S1-Q06, source={Score Advanced - Test 17 (25 May 2023)}, paper=1, qno=6, type=multi, subtopic={Trends}, answer={A,B,C}, answer_src=key, difficulty=5, diff_parts={K2 S1 I1 C1 F1}, diff_note={concepts: statements on periodic properties, trends | steps: compare atomic and ionic properties}, also_in={Score Advanced - Test 35 (15 May 2024) P1 Q6}]
Among the following, the correct statements is/are:
\begin{qoptions}
\opt{A}{The first pK$_a$ of malonic acid is lower than the pK$_a$ of acetic acid while its second pK$_a$ is higher than the pK$_a$ of acetic acid.}
\opt{B}{Picric acid shows positive test with NaHCO$_3$.}
\opt{C}{CH(CN)$_3$ is a stronger acid than HCN.}
\opt{D}{Amides are generally stronger acid than water due to resonance in their conjugate base.}
\end{qoptions}
\end{bankq}
\begin{bankq}[id=SA23-P1-CHEM-S1-Q01, source={Score Advanced - Test 23 (14 Apr 2024)}, paper=1, qno=1, type=single, subtopic={Trends}, answer={D}, answer_src=key, difficulty=4, diff_parts={K2 S1 I1 C1 F0}, diff_note={concepts: correct statement on periodic properties, trends | steps: compare the listed properties}]
Select the correct statements\par
(I) 1 gm - molecule of any gas contains equal number of molecules\par
(II) Atomic mass of carbon-12 is 12u\par
(III) During a reaction, moles and mass of atoms remain constant\par
(IV) Both molality and mole fraction are temperature independent
\begin{qoptions}
\opt{A}{only I, II, III}
\opt{B}{only III, IV}
\opt{C}{only I, III, IV}
\opt{D}{I, II, III and IV}
\end{qoptions}
\end{bankq}
\begin{bankq}[id=SA24-P1-CHEM-S1-Q04, source={Score Advanced - Test 24 (17 Apr 2024)}, paper=1, qno=4, type=multi, subtopic={Periodic properties}, answer={A,B,C,D}, answer_src=key, difficulty=8, diff_parts={K3 S1 I2 C2 F1}, diff_note={concepts: correct statements on periodic properties, compare each property | steps: compare each property}]
Which of the following statements are correct?
\begin{qoptions}
\opt{A}{Linear overlap of two atomic p-orbitals leads to a sigma bond}
\opt{B}{The bond angle H $-$ N $-$ H in NH$_3$ is greater than the bond angle H $-$ As $-$ H in AsH$_3$.}
\opt{C}{Anhydrous HCl is a bad conductor of electricity but aqueous HCl is a good conductor}
\opt{D}{o-nitrophenol is steam volatile whereas p-nitrophenol is not.}
\end{qoptions}
\end{bankq}
\begin{bankq}[id=SA24-P1-CHEM-S3-Q07, source={Score Advanced - Test 24 (17 Apr 2024)}, paper=1, qno=7, type=numerical, subtopic={Periodic properties}, answer={6.1}, answer_src=key, difficulty=8, diff_parts={K3 S1 I2 C2 F1}, diff_note={concepts: effective nuclear charge using Slater rules, periodic properties, shielding constants for each group of electrons | steps: shielding constants for each group of electrons}]
Calculate effective nuclear charge on last electron of Cl. (using Slater rule)
\end{bankq}
\begin{bankq}[id=SA25-P1-CHEM-S3-Q11, source={Score Advanced - Test 25 (21 Apr 2024)}, paper=1, qno=11, type=single, subtopic={Trends}, answer={B}, answer_src=key, difficulty=7, diff_parts={K3 S1 I2 C2 F0}, diff_note={concepts: matching elements with a property, trends, periodic property comparison | steps: periodic property comparison}]
{\footnotesize\setlength{\tabcolsep}{1.5pt}\noindent\begin{tabular}{|c|L{2.6cm}|c|c|} \hline & \multicolumn{1}{c|}{\parbox{2.6cm}{\centering List-I (Element in periodic table)}} & & \multicolumn{1}{c|}{\parbox{2.6cm}{\centering List-II [Graphical change in radius (r) with atomic number (Z)]}}\\ \hline
(P) & Halogens & (1) & \bfigi{59.5}{SA25_p1_chem_s3_q11_G1.png}\\ \hline
(Q) & M$^{+3}$ ions of Lanthanides & (2) & \bfigi{70.3}{SA25_p1_chem_s3_q11_G2.png}\\ \hline
(R) & M atoms of Lanthanides & (3) & \bfigi{61.9}{SA25_p1_chem_s3_q11_G3.png}\\ \hline
(S) & Size of 3d series metal (covalent radius) & (4) & \bfigi{40.5}{SA25_p1_chem_s3_q11_G4.png}\\ \hline
 & & (5) & \bfigi{38.1}{SA25_p1_chem_s3_q11_G5.png}\\ \hline
\end{tabular}}\par
Select \textbf{CORRECT} code :-
\begin{qoptions}
\opt{A}{P $\rightarrow$ 4;Q $\rightarrow$ 5;R $\rightarrow$ 1;S $\rightarrow$ 2}
\opt{B}{P $\rightarrow$ 5;Q $\rightarrow$ 1;R $\rightarrow$ 3;S $\rightarrow$ 2}
\opt{C}{P $\rightarrow$ 4;Q $\rightarrow$ 3;R $\rightarrow$ 1;S $\rightarrow$ 2}
\opt{D}{P $\rightarrow$ 5;Q $\rightarrow$ 1;R $\rightarrow$ 2;S $\rightarrow$ 3}
\end{qoptions}
\end{bankq}
\begin{bankq}[id=SA26-P1-CHEM-S3-Q11, source={Score Advanced - Test 26 (21 Apr 2024)}, paper=1, qno=11, type=single, subtopic={Trends}, answer={B}, answer_src=key, difficulty=7, diff_parts={K3 S1 I2 C2 F0}, diff_note={concepts: matching elements with a property, trends, periodic property comparison | steps: periodic property comparison}]
{\footnotesize\setlength{\tabcolsep}{1.5pt}\noindent\begin{tabular}{|c|L{2.6cm}|c|c|} \hline & \multicolumn{1}{c|}{\parbox{2.6cm}{\centering List-I (Element in periodic table)}} & & \multicolumn{1}{c|}{\parbox{2.6cm}{\centering List-II [Graphical change in radius (r) with atomic number (Z)]}}\\ \hline
(P) & Halogens & (1) & \bfigi{62.4}{SA26_p1_chem_s3_q10_1.png}\\ \hline
(Q) & M$^{+3}$ ions of Lanthanides & (2) & \bfigi{64.3}{SA26_p1_chem_s3_q10_2.png}\\ \hline
(R) & M atoms of Lanthanides & (3) & \bfigi{62.4}{SA26_p1_chem_s3_q10_3.png}\\ \hline
(S) & Size of 3d series metal (covalent radius) & (4) & \bfigi{41.3}{SA26_p1_chem_s3_q10_4.png}\\ \hline
 & & (5) & \bfigi{41.3}{SA26_p1_chem_s3_q10_5.png}\\ \hline
\end{tabular}}\par
Select \textbf{CORRECT} code :-
\begin{qoptions}
\opt{A}{P $\rightarrow$ 4;Q $\rightarrow$ 5;R $\rightarrow$ 1;S $\rightarrow$ 2}
\opt{B}{P $\rightarrow$ 5;Q $\rightarrow$ 1;R $\rightarrow$ 3;S $\rightarrow$ 2}
\opt{C}{P $\rightarrow$ 4;Q $\rightarrow$ 3;R $\rightarrow$ 1;S $\rightarrow$ 2}
\opt{D}{P $\rightarrow$ 5;Q $\rightarrow$ 1;R $\rightarrow$ 2;S $\rightarrow$ 3}
\end{qoptions}
\end{bankq}
\begin{bankq}[id=SA30-P1-CHEM-S2-Q06, source={Score Advanced - Test 30 (04 May 2024)}, paper=1, qno=6, type=multi, subtopic={Trends}, answer={A}, answer_src=key, difficulty=5, diff_parts={K2 S1 I1 C1 F1}, diff_note={concepts: standard problem in this chapter, trends | steps: apply the relevant results; compute}, also_in={Score Advanced - Test 31 (04 May 2024) P1 Q6}]
Which of the following order is/are correct?
\begin{qoptions}
\opt{A}{RbCl $>$ LiCl (Ionic character)}
\opt{B}{AlCl$_3$ $>$ BCl$_3$ (Lewis acidic strength)}
\opt{C}{Atomic nitrogen $>$ molecular nitrogen (Ionization energy)}
\opt{D}{Li$_2$CO$_3$ $>$ Cs$_2$CO$_3$ (Thermal stability)}
\end{qoptions}
\end{bankq}
\begin{bankq}[id=SA33-P2-CHEM-S2-Q09, source={Score Advanced - Test 33 (12 May 2024)}, paper=2, qno=9, type=multi, subtopic={Trends}, answer={A,B,C}, answer_src=key, difficulty=8, diff_parts={K3 S1 I2 C2 F1}, diff_note={concepts: standard problem in this chapter, trends, apply the relevant results | steps: apply the relevant results; compute}, also_in={Score Advanced - Test 34 (12 May 2024) P2 Q9}]
Select the correct comparison(s) of properties indicated in options :
\begin{qoptions}
\opt{A}{MnO$_4^{2-}$ $>$ MnO$_4^-$ : (paramagnetic character)}
\opt{B}{Na$_2$Cr$_2$O$_7$ $>$ K$_2$Cr$_2$O$_7$ (Solubility in water)}
\opt{C}{Lu$^{+3}$ $<$ Ce$^{+3}$ $<$ La$^{+3}$ : (Ionic radius)}
\opt{D}{Mn$^{+3}$ $<$ Fe$^{+3}$ $<$ Ti$^{+3}$ : (oxidising power)}
\end{qoptions}
\end{bankq}

\section*{SRGBT papers}
\begin{bankq}[id=SRGBT03-P1-CHEM-S1-Q21, source={SRGBT - Test 03 (08 Aug 2021)}, paper=1, qno=21, type=multi, subtopic={Trends}, answer={B,C,D}, answer_src=key, difficulty=8, diff_parts={K3 S1 I2 C2 F1}, diff_note={concepts: standard problem in this chapter, trends, apply the relevant results | steps: apply the relevant results; compute}]
Select incorrect graphs, which are plotted between ionisation energy of 2nd period elements and their atomic number.\par
\bfig{294.8}{SRGBT03_p1_chem_s0_q21_1.png}
\end{bankq}
\begin{bankq}[id=SRGBT03-P1-CHEM-S1-Q23, source={SRGBT - Test 03 (08 Aug 2021)}, paper=1, qno=23, type=multi, subtopic={Trends}, answer={A,B,C}, answer_src=key, difficulty=8, diff_parts={K3 S1 I2 C2 F1}, diff_note={concepts: standard problem in this chapter, trends, apply the relevant results | steps: apply the relevant results; compute}]
Choose the incorrect ionisation energy order.
\begin{qoptions}
\opt{A}{Li $<$ B $<$ Be $<$ C $<$ N $<$ O $<$ F}
\opt{B}{Na $<$ Mg $<$ Al $<$ Si $<$ S $<$ P $<$ Cl}
\opt{C}{K $<$ Ga $<$ Ca $<$ Ge $<$ As $<$ Se $<$ Br}
\opt{D}{Rb $<$ Sr $<$ In $<$ Sn $<$ Sb $<$ Te $<$ I}
\end{qoptions}
\end{bankq}
\begin{bankq}[id=SRGBT03-P1-CHEM-S1-Q25, source={SRGBT - Test 03 (08 Aug 2021)}, paper=1, qno=25, type=multi, subtopic={Trends}, answer={A,B,D}, answer_src=key, difficulty=8, diff_parts={K3 S1 I2 C2 F1}, diff_note={concepts: standard problem in this chapter, trends, apply the relevant results | steps: apply the relevant results; compute}]
In which of the following arrangement, the order is according to the property indicated against it:
\begin{qoptions}
\opt{A}{I $<$ Br $<$ F $<$ Cl : Increasing EA}
\opt{B}{Li $<$ Na $<$ K $<$ Rb : Increasing metallic radius}
\opt{C}{B $<$ C $<$ N $<$ O : Increasing frist ionization energy}
\opt{D}{Al$^{+3}$ $<$ Mg$^{+2}$ $<$ Na$^+$ $<$ F$^-$ : Increasing ionic size}
\end{qoptions}
\end{bankq}
\begin{bankq}[id=SRGBT03-P1-CHEM-S1-Q27, source={SRGBT - Test 03 (08 Aug 2021)}, paper=1, qno=27, type=multi, subtopic={Trends}, answer={C,D}, answer_src=key, difficulty=8, diff_parts={K3 S1 I2 C2 F1}, diff_note={concepts: standard problem in this chapter, trends, apply the relevant results | steps: apply the relevant results; compute}]
Following is the graph representing Electronegativity of certain elements A$_1$, A$_2$, A$_3$ , A$_4$, A$_5$. Which of the following option(s) representing corresponding elements can be \textbf{correct}?\par
\bfig{102.0}{SRGBT03_p1_chem_s0_q27_1.png}
\begin{qoptions}
\opt{A}{A$_1$ = \textbf{N}, A$_2$ = \textbf{P}, A$_3$ = \textbf{F}, A$_4$ = \textbf{Cl}, A$_5$ = \textbf{Al}}
\opt{B}{A$_1$ = \textbf{Cl}, A$_2$ = \textbf{P}, A$_3$ = \textbf{O}, A$_4$ = \textbf{N}, A$_5$ = \textbf{Al}}
\opt{C}{A$_1$ = \textbf{N}, A$_2$ = \textbf{P}, A$_3$ = \textbf{O}, A$_4$ = \textbf{Si}, A$_5$ = \textbf{Al}}
\opt{D}{A$_1$ = \textbf{N}, A$_2$ = \textbf{P}, A$_3$ = \textbf{F}, A$_4$ = \textbf{Al}, A$_5$ = \textbf{Sr}}
\end{qoptions}
\end{bankq}
\begin{bankq}[id=SRGBT03-P1-CHEM-S1-Q29, source={SRGBT - Test 03 (08 Aug 2021)}, paper=1, qno=29, type=multi, subtopic={Periodic properties}, answer={B,C}, answer_src=key, difficulty=8, diff_parts={K3 S1 I2 C2 F1}, diff_note={concepts: standard problem in this chapter, periodic properties, apply the relevant results | steps: apply the relevant results; compute}]
Which of the following statements is/are not correct.
\begin{qoptions}
\opt{A}{Na$_2$O $<$ MgO $<$ ZnO $<$ P$_4$O$_{10}$ : Acidic property.}
\opt{B}{Na $<$ Si $>$ Mg $<$ Al : First ionisation potential.}
\opt{C}{F $>$ Cl $>$ Br : Electron affinity.}
\opt{D}{H$^-$ $>$ F$^-$ : Size}
\end{qoptions}
\end{bankq}
\begin{bankq}[id=SRGBT03-P1-CHEM-S1-Q32, source={SRGBT - Test 03 (08 Aug 2021)}, paper=1, qno=32, type=multi, subtopic={Periodic properties}, answer={130.56}, answer_src=key, difficulty=8, diff_parts={K3 S1 I2 C2 F1}, diff_note={concepts: standard problem in this chapter, periodic properties, apply the relevant results | steps: apply the relevant results; compute}]
Calculate the 10$^\mathrm{th}$ ionisation energy of Neon in Mega Joules/mole.\par
[Use N$_\mathrm A$ = 6 $\times$ 10$^{23}$ \& 1 eV = 1.6 $\times$ 10$^{-19}$ J]
\end{bankq}
\begin{bankq}[id=SRGBT03-P1-CHEM-S1-Q36, source={SRGBT - Test 03 (08 Aug 2021)}, paper=1, qno=36, type=multi, subtopic={Periodic properties}, answer={728}, answer_src=key, difficulty=8, diff_parts={K3 S1 I2 C2 F1}, diff_note={concepts: standard problem in this chapter, periodic properties, apply the relevant results | steps: apply the relevant results; compute}]
When electrons (one from each atom) are removed from 0.5 mol of A and given them (one to each atom) to 0.5 mol of B, 40 kJ of energy is absorbed. When electrons (one from each atom) are removed from 0.5 mol of B and given them (one to each atom) to 0.5 mol of A, 52 kJ energy is released. `When electrons (two from each atom) are removed from 0.5 mol of A and given them (one to each atom) to one mole of B, 280 kJ energy is absorbed\par
IE$_1$ of A = 400 KJ/mol\par
Electron gain enthalpy of A = $-$352 KJ/mol\par
Sample of A and B are taken in gas phase.\par
(i) Ionisation energy of B (in KJ/mol) = `X'\par
(ii) Second ionisation energy of A (in KJ/mol) = ``Y''\par
(iii) Electron gain enthalpy of B (in KJ/mol) = ``Z''\par
\textbf{Then find the value of (X + Y + Z).}
\end{bankq}
\begin{bankq}[id=SRGBT03-P1-CHEM-S1-Q40, source={SRGBT - Test 03 (08 Aug 2021)}, paper=1, qno=40, type=multi, subtopic={Periodic properties}, answer={2}, answer_src=key, difficulty=8, diff_parts={K3 S1 I2 C2 F1}, diff_note={concepts: standard problem in this chapter, periodic properties, apply the relevant results | steps: apply the relevant results; compute}]
IE$_1$ (kJ/mol) of 2nd period elements are mentioned below with some blank/vacant data spaces.\par
{\footnotesize\setlength{\tabcolsep}{2pt}\noindent\begin{tabular}{cccccccc}
Li & Be & B & C & N & O & F & Ne \\
520 & \underline{``P''} & 800 & \underline{``Q''} & 1402 & 1314 & \underline{``R''} & 2080 \\
\end{tabular}}\par
Select correct number of statement(s) regarding above data table. (Consider all processes are gas phased)\par
(i) $|$EA$|$ of N$^+$ is more than 1402 kJ/mol\par
(ii) Value of ``R'' is more than 1402 kJ/mol\par
(iii) Value of P is more than $|$EA$|$ of C$^+$\par
(iv) Q $>$ 800 kJ/mol\par
(v) $\dfrac12$ mole O$^+\to$ O released energy = S ; $\dfrac12$ mole Ne$^+\to$ Ne released energy = T, then ($|$S$|$ + $|$T$|$ $>$ 2R) is correct.)
\end{bankq}
\begin{bankq}[id=SRGBT10-P1-CHEM-S1-Q22, source={SRGBT - Test 10 (26 Sep 2021)}, paper=1, qno=22, type=multi, subtopic={Trends}, answer={A,B}, answer_src=key, difficulty=8, diff_parts={K3 S1 I2 C2 F1}, diff_note={concepts: standard problem in this chapter, trends, apply the relevant results | steps: apply the relevant results; compute}]
Select the option(s) having correct order of the properties involved.
\begin{qoptions}
\opt{A}{Size in gas phase : Al$^{3+}$ $<$ Ga$^{3+}$ $<$ H$^{-}$}
\opt{B}{IE$_1$ values : Sb $<$ Te}
\opt{C}{At. size : Ag $>$ Hg}
\opt{D}{IE$_1$ values : Bi $>$ Sb}
\end{qoptions}
\end{bankq}
\begin{bankq}[id=SRGBT10-P1-CHEM-S1-Q24, source={SRGBT - Test 10 (26 Sep 2021)}, paper=1, qno=24, type=multi, subtopic={Periodic properties}, answer={A,B,C}, answer_src=key, difficulty=8, diff_parts={K3 S1 I2 C2 F1}, diff_note={concepts: standard problem in this chapter, periodic properties, apply the relevant results | steps: apply the relevant results; compute}]
One mole of `F' atoms are ionised to F$^-$ the energy released is `X' Joules. Then the incorrect statements:
\begin{qoptions}
\opt{A}{X Joules is sufficient to ionise 1 mole of gaseous Cl into Cl$^+$.}
\opt{B}{X Joules is sufficient to ionise 1 mole of gaseous F into F$^+$.}
\opt{C}{X Joules is sufficient to ionise 1 mole of gaseous F into F$^+$ as well as 1 mole of gaseous Cl into Cl$^+$.}
\opt{D}{Less than 1 mole of gaseous F or Cl will be ionised to F$^+$ or Cl$^+$.}
\end{qoptions}
\end{bankq}
\begin{bankq}[id=SRGBT10-P1-CHEM-S1-Q26, source={SRGBT - Test 10 (26 Sep 2021)}, paper=1, qno=26, type=multi, subtopic={Periodic properties}, answer={A,C}, answer_src=key, difficulty=8, diff_parts={K3 S1 I2 C2 F1}, diff_note={concepts: standard problem in this chapter, periodic properties, apply the relevant results | steps: apply the relevant results; compute}]
X, Y and Z are three consecutive elements of second period whose I.P. data (in eV) are available. By using these data specify incorrect statements regarding X, Y and Z.\par
\begin{tabular}{|c|c|c|c|c|}\hline
Element & IP$_1$ & IP$_2$ & IP$_3$ & IP$_4$\\\hline
X & 5.39 & 75.6 & 122.4 & Not possible\\\hline
Y & 8.2 & 25.1 & 37.9 & 259.2\\\hline
Z & 9.3 & 18.2 & 217.6 & 363.7\\\hline
\end{tabular}\par
\begin{qoptions}
\opt{A}{X$_{(g)}$ + e$^-$ $\longrightarrow$ A$^-_{(g)}$ \quad $\Delta$H = positive}
\opt{B}{Z$_\mathrm Y$ $>$ Z$_\mathrm Z$ (Z is atomic number)}
\opt{C}{Oxide of `Y' is amphoteric}
\opt{D}{X$^{+1}$ is stable than X$^{2+}$}
\end{qoptions}
\end{bankq}
\begin{bankq}[id=SRGBT10-P1-CHEM-S1-Q28, source={SRGBT - Test 10 (26 Sep 2021)}, paper=1, qno=28, type=multi, subtopic={Trends}, answer={A,B,C}, answer_src=key, difficulty=8, diff_parts={K3 S1 I2 C2 F1}, diff_note={concepts: standard problem in this chapter, trends, apply the relevant results | steps: apply the relevant results; compute}]
Select the option containing correct order of the parameters given
\begin{qoptions}
\opt{A}{Increasing order of acidic strength : Na$_2$O, MgO, CO$_2$, SO$_3$}
\opt{B}{Decreasing order of ionic radii : K$^+$, Na$^+$, Li$^+$, Mg$^{2+}$}
\opt{C}{Decreasing order of boiling points : H$_2$O, HF, NH$_3$}
\opt{D}{Decreasing order of boiling points : He, Ne, Ar, Xe}
\end{qoptions}
\end{bankq}
\begin{bankq}[id=SRGBT10-P1-CHEM-S1-Q29, source={SRGBT - Test 10 (26 Sep 2021)}, paper=1, qno=29, type=multi, subtopic={Periodic properties}, answer={B,C,D}, answer_src=key, difficulty=8, diff_parts={K3 S1 I2 C2 F1}, diff_note={concepts: standard problem in this chapter, periodic properties, apply the relevant results | steps: apply the relevant results; compute}]
X and Y are second period elements. They form two compounds XF$_3$ and YH$_3$ respectively. XF$_3$ is resonance stabilized while YH$_3$ is not. Select the incorrect statement.
\begin{qoptions}
\opt{A}{X--F bond distance in XF$_3$ is smaller than that in XF$_4^-$.}
\opt{B}{X--F bond in XF$_3$ is larger than that in XF$_3^-$ -- YH$_3^+$.}
\opt{C}{XF$_3$ does not have back bonding.}
\opt{D}{YH$_3$ is a Lewis acid.}
\end{qoptions}
\end{bankq}
\begin{bankq}[id=SRGBT10-P1-CHEM-S1-Q30, source={SRGBT - Test 10 (26 Sep 2021)}, paper=1, qno=30, type=multi, subtopic={Trends}, answer={A,B,C,D}, answer_src=key, difficulty=8, diff_parts={K3 S1 I2 C2 F1}, diff_note={concepts: standard problem in this chapter, trends, apply the relevant results | steps: apply the relevant results; compute}]
Which of the following statements are \textbf{correct} :
\begin{qoptions}
\opt{A}{IE$_2$ of Cu $>$ IE$_2$ of Zn where as IE$_1$ of Cu $<$ IE$_1$ of Zn.}
\opt{B}{$|\Delta$H$_\mathrm{EG}|$ of S $>$ $|\Delta$H$_\mathrm{EG}|$ of O.}
\opt{C}{Reducing nature of Li in aqueous solution is more than that of Na.}
\opt{D}{There is a substantial contraction in size observed on moving from 3$^\mathrm{rd}$ group to 4th group in 6th period of periodic table.}
\end{qoptions}
\end{bankq}
\begin{bankq}[id=SRGBT10-P1-CHEM-S1-Q32, source={SRGBT - Test 10 (26 Sep 2021)}, paper=1, qno=32, type=multi, subtopic={Trends}, answer={4}, answer_src=key, difficulty=8, diff_parts={K3 S1 I2 C2 F1}, diff_note={concepts: standard problem in this chapter, trends, apply the relevant results | steps: apply the relevant results; compute}, also_in={Nurture (General) - Test 01 (12 Jun 2021) P1 Q2}]
How many comparison(s) are \textbf{INCORRECT} among given ?\par
(i) H $>$ Li (ionisation energy)\quad (ii) Li $>$ Be (size)\quad (iii) Na $>$ Rb (size)\quad (iv) O $>$ N (size)\quad (v) S $>$ P (ionisation energy)\quad (vi) C $>$ O (size)\quad (vii) B $>$ Al (size)\quad (viii) F $>$ Cl (ionisation energy)\quad (ix) N $>$ O (ionisation energy)
\end{bankq}
\begin{bankq}[id=SRGBT10-P1-CHEM-S1-Q34, source={SRGBT - Test 10 (26 Sep 2021)}, paper=1, qno=34, type=multi, subtopic={Periodic properties}, answer={6.7}, answer_src=key, difficulty=8, diff_parts={K3 S1 I2 C2 F1}, diff_note={concepts: standard problem in this chapter, periodic properties, apply the relevant results | steps: apply the relevant results; compute}]
Assuming screening by K shell electrons to be same in Li \& Be. Calculate I.E.$_1$ (in eV) of Li if IE$_2$ of Be is 13.6$\times$1.44 eV/atom.
\end{bankq}
\begin{bankq}[id=SRGBT10-P1-CHEM-S1-Q38, source={SRGBT - Test 10 (26 Sep 2021)}, paper=1, qno=38, type=multi, subtopic={Trends}, answer={5}, answer_src=key, difficulty=8, diff_parts={K3 S1 I2 C2 F1}, diff_note={concepts: standard problem in this chapter, trends, apply the relevant results | steps: apply the relevant results; compute}]
The number of correct trends is/are :\par
(i) AgF $>$ AgCl $>$ AgBr $>$ AgI (Solubility in water)\par
(ii) TiO$_2$ $<$ SiO$_2$ (Polarising power)\par
(iii) O$_2^-$ $>$ O$_2^{2-}$ (Ionic radii)\par
(iv) Cu $>$ Ag $<$ Au (Ionisation energy)\par
(v) N $<$ O $<$ F $<$ Ne (Z$_\mathrm{eff}$)\par
(vi) La$^{3+}$ $>$ Nd$^{3+}$ $>$ Gd$^{3+}$ $>$ Ho$^{3+}$ $>$ Lu$^{3+}$ (Ionic radii)
\end{bankq}

\section*{TOAS papers}
\begin{bankq}[id=TOAS01-P2-CHEM-S2-Q05, source={TOAS - Test 01 (30 Jan 2022)}, paper=2, qno=5, type=numerical, subtopic={Periodic properties}, answer={2}, answer_src=key, difficulty=3, diff_parts={K1 S2 I0 C0 F1}, diff_note={concepts: atomic and ionic sizes | steps: group trends (O, S) and (Li, Na); noble gas radii (He, Ne), (Kr, Ne); cation, anion comparisons}]
Write the number of pairs in which size of first element or ion is higher as compare to II$^{\mathrm{nd}}$ out of following pairs :\\
(O, S), (He, Ne), (Kr, Ne), ($\mathrm{Na}$, $\mathrm{Na^+}$), (Cl, $\mathrm{Cl^-}$), (Li, Na), ($\mathrm{Li^+}$, $\mathrm{Na^+}$)
\end{bankq}
\begin{bankq}[id=TOAS01-P2-CHEM-S2-Q06, source={TOAS - Test 01 (30 Jan 2022)}, paper=2, qno=6, type=numerical, subtopic={Periodic properties}, answer={5}, answer_src=solved, working={First IE (eV): Be 9.3; Mg 7.6, Ca 6.1, Sr 5.7, Ba 5.2, Li 5.4 - all five are lower}, difficulty=2, diff_parts={K1 S1 I0 C0 F1}, diff_note={concepts: ionisation energies | steps: Be is higher than Mg, Ca, Sr, Ba and Li}]
Find the number of element having first ionization energy lower than Be.\\
Mg, Ca, Sr, Ba, Li
\end{bankq}
\begin{bankq}[id=TOAS01-P2-CHEM-S2-Q07, source={TOAS - Test 01 (30 Jan 2022)}, paper=2, qno=7, type=numerical, subtopic={Periodic properties}, answer={3}, answer_src=solved, working={(a) false: HOCl>HOBr>HOI; (b) false: BeO is amphoteric; (c), (d), (e) true}, difficulty=3, diff_parts={K1 S2 I0 C0 F1}, diff_note={concepts: periodic trends | steps: acidic strength and basic strength; electronegativity and electron affinity; ionic character}]
How many orders are CORRECT\\
(a) Acidic strength $\Rightarrow$ ClOH $<$ BrOH $<$ IOH\\
(b) Basic strength $\Rightarrow$ MgO $<$ CaO $<$ SrO $<$ BeO\\
(c) Electronegativity $\Rightarrow$ I $<$ Br $<$ N $<$ O $<$ F\\
(d) Electron affinity $\Rightarrow$ Mg $<$ Na $<$ Si $<$ S $<$ Cl\\
(e) \% Ionic character $\Rightarrow$ NaCl $<$ KF $<$ RbF
\end{bankq}
\begin{bankq}[id=TOAS05-P1-CHEM-S1-Q08, source={TOAS - Test 05 (03 Apr 2022)}, paper=1, qno=8, type=multi, subtopic={Periodic properties}, answer={A,C}, answer_src=key, difficulty=2, diff_parts={K1 S1 I0 C0 F1}, diff_note={concepts: periodic trends | steps: identify the property in each option; compare the order}]
Choose the INCORRECT order of the indicated property given in the options.
{\small\setlength{\tabcolsep}{3pt}\renewcommand{\arraystretch}{1.25}\begin{center}\begin{tabular}{|L{0.6cm}|L{3.3cm}|c|L{2.3cm}|}
\hline
(A) & $\mathrm{He<Ne<Kr<Ar}$ & : & Clathrate compound formation tendency order\\
\hline
(B) & $\mathrm{NaO_2}<\mathrm{KO_2}<\mathrm{RbO_2}<\mathrm{CsO_2}$ & : & Thermal stability order\\
\hline
(C) & $\mathrm{BeF_2<BaF_2}$ & : & Solubility order\\
\hline
(D) & $\mathrm{LiF>LiCl>LiI}$ & : & Melting point order\\
\hline
\end{tabular}\end{center}}
\end{bankq}
\begin{bankq}[id=TOAS09-P1-CHEM-S2-Q05, source={TOAS - Test 09 (30 May 2022)}, paper=1, qno=5, type=multi, subtopic={Periodic properties}, answer={A,B,C,D}, answer_src=key, difficulty=2, diff_parts={K1 S1 I0 C0 F1}, diff_note={concepts: atomic size trends | steps: lanthanide contraction and transition series trends; Cu and Zn sizes}]
The CORRECT order(s) regarding the atomic size is / are
\begin{qoptions}
\opt{A}{Zr $\simeq$ Hf}
\opt{B}{Fe $\simeq$ Co $\simeq$ Ni}
\opt{C}{Sc $<$ Y $<$ La}
\opt{D}{Cu $<$ Zn}
\end{qoptions}
\end{bankq}
\begin{bankq}[id=TOAS09-P2-CHEM-S1-Q01, source={TOAS - Test 09 (30 May 2022)}, paper=2, qno=1, type=single, subtopic={Periodic properties}, answer={D}, answer_src=key, difficulty=1, diff_parts={K1 S1 I0 C0 F0}, diff_note={concepts: second ionisation energies | steps: electronic configurations of the ions; compare the second IE order}]
The correct order of second I.E. of C, N, O and F are in the order
\begin{qoptions}
\opt{A}{F $>$ O $>$ N $>$ C}
\opt{B}{C $>$ N $>$ O $>$ F}
\opt{C}{O $>$ N $>$ F $>$ C}
\opt{D}{O $>$ F $>$ N $>$ C}
\end{qoptions}
\end{bankq}
\begin{bankq}[id=TOAS13-P2-CHEM-S1-Q14, source={TOAS - Test 13 (24 Jul 2022)}, paper=2, qno=14, type=single, subtopic={Trends}, answer={C}, answer_src=key, difficulty=1, diff_parts={K1 S1 I0 C0 F0}, diff_note={concepts: periodic and group trends | steps: polarising power, thermal stability and bond order orders are reversed; Lewis acidity of the boron halides rises down the group}]
Select the \textbf{CORRECT} option, for mentioned properties
\begin{qoptions}
\opt{A}{Ba$^{2+}$ $>$ Ca$^{2+}$ $>$ Mg$^{2+}$ $>$ Be$^{2+}$ : (polarising power of cation)}
\opt{B}{BeSO$_4$ $>$ MgSO$_4$ $>$ CaSO$_4$ $>$ SrSO$_4$ : (thermal stability)}
\opt{C}{BF$_3$ $<$ BCl$_3$ $<$ BBr$_3$ $<$ BI$_3$ : (lewis acidic nature)}
\opt{D}{O$_2$ $<$ KO$_2$ $<$ K$_2$O$_2$ : (O--O bond order)}
\end{qoptions}
\end{bankq}
\begin{bankq}[id=TOAS17-P2-CHEM-S1-Q16, source={TOAS - Test 17 (18 Sep 2022)}, paper=2, qno=16, type=single, subtopic={Periodic properties}, answer={A}, answer_src=key, difficulty=2, diff_parts={K1 S2 I0 C0 F0}, diff_note={concepts: periodic properties | steps: reducing power; softness; stability; acidity}]
Choose incorrect order regarding mentioned properties :
\begin{qoptions}
\opt{A}{SeO$_2$ $>$ SO$_2$ (Reducing power)}
\opt{B}{Na $>$ K (Softness of metal)}
\opt{C}{PbCl$_4$ $>$ PbI$_4$ (stability)}
\opt{D}{H$_2$S $>$ H$_2$O (Acidic character)}
\end{qoptions}
\end{bankq}
\begin{bankq}[id=TOAS19-P1-CHEM-S1-Q03, source={TOAS - Test 19 (05 Feb 2023)}, paper=1, qno=3, type=multi, subtopic={Trends}, answer={A,B,C}, answer_src=key, difficulty=3, diff_parts={K1 S2 I0 C0 F1}, diff_note={concepts: periodic trends | steps: solubility of the chlorides; hydride stability; back donation in silyl amines; acidity of silanol}]
Select the correct order(s) with respect to specified property.
\begin{qoptions}
\opt{A}{NaCl $>$ KCl (solubility in water)}
\opt{B}{LiH $>$ NaH (Thermal stability)}
\opt{C}{N(CH$_3$)$_3$ $>$ N(SiH$_3$)$_3$ (lewis basic strength)}
\opt{D}{(CH$_3$)$_3$SiOH $<$ (CH$_3$)$_3$COH (H$^+$ donating ability)}
\end{qoptions}
\end{bankq}

\section*{TOPS papers}
\begin{bankq}[id=TOPS01-P1-CHEM-S3-Q03, source={TOPS - Test 01 (31 Jul 2022)}, paper=1, qno=3, type=numerical, subtopic={Periodic properties}, answer={6.00}, answer_src=key, difficulty=6, diff_parts={K2 S2 I1 C1 F1}, diff_note={concepts: periodic trends of oxides, oxidation state | steps: identify the period and A; most acidic oxide; least acidic oxide; difference in the oxidation states}]
Element ``A'' only forms acidic oxides. ``A'' belongs to the period containing Bridging elements.

If A exists in nature as A$_2$ ; then the difference in oxidation states of A in most acidic \& least acidic oxides of A is \underline{\hspace{1cm}}.
\end{bankq}
\begin{bankq}[id=TOPS01-P1-CHEM-S3-Q06, source={TOPS - Test 01 (31 Jul 2022)}, paper=1, qno=6, type=numerical, subtopic={Periodic properties}, answer={4.00}, answer_src=key, difficulty=7, diff_parts={K3 S2 I1 C1 F1}, diff_note={concepts: electron affinity, ionisation energy, bond angle | steps: check each of seven orders; ionisation energy irregularities; electron count | seven small checks}]
Number of correct order(s) among the following

(i) Se $>$ O (EA)

(ii) Na $>$ Al (EA)

(iii) H$_2$O $>$ NH$_3$ (Bond Angle)

(iv) Pd $>$ Ni (IE)

(v) N$_3^-$ $>$ CO$_2$ (number of electrons)

(vi) Mn $>$ Fe (IE$_3$)

(vii) B $>$ H (EN)
\end{bankq}
\begin{bankq}[id=TOPS03-P1-CHEM-S2-Q10, source={TOPS - Test 03 (20 Nov 2022)}, paper=1, qno=10, type=multi, subtopic={Periodic properties}, answer={A,B,D}, answer_src=key, difficulty=7, diff_parts={K3 S2 I1 C1 F1}, diff_note={concepts: ionic size, ionisation energy, electronegativity | steps: isoelectronic sizes; second ionisation energy irregularities; electronegativity; metallic radius | four trends}]
In which of the following arrangements. The order is correct according to the property indicating against it :-
\begin{qoptions}
\opt{A}{Increasing size : Al$^{3+}<$ Mg$^{+2}<$ Na$^+<$ F$^\Theta$}
\opt{B}{Increasing IE$_2$ : C $<$ B $<$ N $<$ O}
\opt{C}{Increasing EN : I $<$ Br $<$ F $<$ Cl}
\opt{D}{Increasing metallic radius : Li $<$ Na $<$ K $<$ Rb}
\end{qoptions}
\end{bankq}
\begin{bankq}[id=TOPS03-P2-CHEM-S3-Q05, source={TOPS - Test 03 (20 Nov 2022)}, paper=2, qno=5, type=numerical, subtopic={Periodic properties}, answer={5}, answer_src=key, difficulty=7, diff_parts={K3 S2 I1 C1 F1}, diff_note={concepts: periodic trends, ionic radius, ionisation energy | steps: check each of six orders; compare electron affinities; compare radii and ionisation energies; count | six orders}]
Find the total number of correct orders among the following.

(1) Order of electron affinity : S $>$ O$^-$

(2) Order of ionic radius : Al$^{3+}<$ Mg$^{2+}<$ O$^{2-}$

(3) Order of Electronegativity : Si $<$ P $<$ C $<$ N

(4) Order of atomic radius : O $<$ S $<$ Se

(5) Order of second ionization energy : C $<$ N $<$ F $<$ O

(6) Order of ionic radii : S$^{2-}<$ Cl$^-<$ K$^+<$ Ca$^{2+}$
\end{bankq}
\begin{bankq}[id=TOPS03-P2-CHEM-S3-Q06, source={TOPS - Test 03 (20 Nov 2022)}, paper=2, qno=6, type=numerical, subtopic={Periodic properties}, answer={7}, answer_src=key, difficulty=5, diff_parts={K1 S2 I1 C1 F1}, diff_note={concepts: acidic basic and amphoteric oxides | steps: nature of each oxide; which react with alkali; count | eight oxides}]
How many of the following can react with alkali :

Mn$_2$O$_7$, ZnO, Al$_2$O$_3$, BeO, SeO$_2$, Cl$_2$O$_7$, Fe$_2$O$_3$, As$_2$O$_3$
\end{bankq}
\begin{bankq}[id=TOPS04-P1-CHEM-S1-Q05, source={TOPS - Test 04 (18 Dec 2022)}, paper=1, qno=5, type=multi, subtopic={Periodic properties}, answer={B,C}, answer_src=key, difficulty=7, diff_parts={K2 S2 I2 C1 F1}, diff_note={concepts: periodic properties, ionisation energy | steps: identify possible elements; compare ionisation per atom; size and shape cases}]
If $A_0$ atoms of element $X$ have ionization energy $E_1$ where number of ``$s$'' electrons in $X$ is less than 5 and $B_0$ atoms of element $Y$ have ionization energy $E_2$ where number of `$s$' electrons in $Y$ is less than 4. Then.
\begin{qoptions}
\opt{A}{$\dfrac{E_2}{B_0}$ must be greater than $\dfrac{E_1}{A_0}$}
\opt{B}{Size of $Y$ may be greater than $X$ if $Y$ belongs to group 1 \& $X$ belongs to group 16}
\opt{C}{$\dfrac{E_2}{B_0}$ will be greater than $\dfrac{E_1}{A_0}$ if $Y$ belongs to group 18}
\opt{D}{If $Y$ exists in nature as $Y_2$ \& $X$ exists in nature $X_2$ then the compound $X_2Y_2$ will always be planar}
\end{qoptions}
\end{bankq}
\begin{bankq}[id=TOPS04-P2-CHEM-S1-Q17, source={TOPS - Test 04 (18 Dec 2022)}, paper=2, qno=17, type=single, subtopic={Periodic properties}, answer={C}, answer_src=key, difficulty=6, diff_parts={K3 S2 I1 C1 F0}, diff_note={concepts: solubility, thermal stability, polarisation | steps: Fajans rule application; trends; check each}]
Which of the following order is correct?
\begin{qoptions}
\opt{A}{AgCl $<$ AgBr (Solubility in water)}
\opt{B}{$\mathrm{BeCO_3>BaCO_3}$ (Thermal Stability)}
\opt{C}{$\mathrm{K^+<Cu^+}$ (Polarizing power)}
\opt{D}{KF $>$ NaF (Covalent character)}
\end{qoptions}
\end{bankq}
\begin{bankq}[id=TOPS04-P2-CHEM-S2-Q09, source={TOPS - Test 04 (18 Dec 2022)}, paper=2, qno=9, type=numerical, subtopic={Periodic properties}, answer={11}, answer_src=solved, working={A is Cl (one unpaired electron, high electron affinity), B is I, compound $\mathrm{ICl_3}$: 2 lone pairs on I plus 3 on each Cl gives 11}, difficulty=6, diff_parts={K2 S2 I1 C1 F1}, diff_note={concepts: periodic trends, structure | steps: identify A and B; structure of the compound; lone pairs}]
The spin only magnetic moment of an atom of element ``$A$'' in its ground state is 1.73 BM. This element belongs to p-block and has significantly high EA value. Another element ``$B$'' of same group and having less EA than $A$ forms a compound $\mathrm{BA_3}$. Elemental form of $A$ in nature is $\mathrm{A_2}$ (gas) \& that of $B$ is $\mathrm{B_2}$ (solid). Then number of lone pairs in $\mathrm{BA_3}$ is \underline{\hspace{1cm}}.
\end{bankq}
\begin{bankq}[id=TOPS05-P2-CHEM-S3-Q06, source={TOPS - Test 05 (14 Jan 2023)}, paper=2, qno=6, type=numerical, subtopic={Periodic properties}, answer={1}, answer_src=key, difficulty=7, diff_parts={K3 S2 I1 C1 F1}, diff_note={concepts: periodic trends, solubility, conductivity | steps: check each of seven orders; count the correct; count the incorrect; difference}]
$X=$ Number of incorrect statement(s) among the following

$Y=$ Number of correct statement(s) among the following

Find $Y-X$ :

(i) Covalent character order : $\mathrm{ZnCl_2<CdCl_2<HgCl_2}$

(ii) Acidic nature of oxide : $\mathrm{CrO_3>CrO>Cr_2O_3}$

(iii) Solubility in water : $\mathrm{BeSO_4>MgSO_4>CaSO_4}$ $\mathrm{>SrSO_4>BaSO_4}$

(iv) Electrical conductivity in infinitely aqueous dilute solution : $\mathrm{CsCl>NaCl}$

(v) Thermal stability : $\mathrm{LiOH<NaOH<KOH}$ $\mathrm{<RbOH<CsOH}$

(vi) Intensity of colour : $\mathrm{AgCl>AgI}$

(vii) Solubility in acetone : $\mathrm{NaCl>NaI}$
\end{bankq}
\begin{bankq}[id=TOPS05-P2-CHEM-S3-Q08, source={TOPS - Test 05 (14 Jan 2023)}, paper=2, qno=8, type=numerical, subtopic={Periodic properties}, answer={1}, answer_src=key, difficulty=8, diff_parts={K3 S2 I2 C1 F1}, diff_note={concepts: electron affinity, ionisation energy, magnetism | steps: data for oxygen; compare each species on the three properties; count}]
Number of species among the following having less EA, more IE and more magnetic moment than O-atom.

B, Al, Si, F, Cl, $\mathrm{O^-}$, N, C, S, P.
\end{bankq}

\section*{TNAS papers}
\begin{bankq}[id=TNAS02-P2-CHEM-S1-Q04, source={TNAS - Test 02 (07 May 2023)}, paper=2, qno=4, type=multi, subtopic={Periodic properties}, answer={C,D}, answer_src=key, difficulty=6, diff_parts={K1 S2 I2 C1 F1}, diff_note={concepts: periodic and ionic character trends | steps: each statement; Fajans rule; solubility}]
Which of the following set of order is \textbf{INCORRECT} for their indicated properties?
\begin{qoptions}
\opt{A}{Covalent character : $\mathrm{CuCl<AgCl<AuCl}$}
\opt{B}{Stability : $\mathrm{PbF_4>PbCl_4>PbBr_4>PbI_4}$}
\opt{C}{Acidic nature order : $\mathrm{BeO<MgO<CaO<SrO}$}
\opt{D}{Water solubility : {\small$\mathrm{LiNO_3<NaNO_3<KNO_3<RbNO_3}$}}
\end{qoptions}
\end{bankq}
\begin{bankq}[id=TNAS03-P1-CHEM-S1-Q06, source={TNAS - Test 03 (28 May 2023)}, paper=1, qno=6, type=single, subtopic={Periodic properties}, answer={D}, answer_src=key, difficulty=3, diff_parts={K1 S1 I1 C1 F0}, diff_note={concepts: ionisation energy and bond order | steps: each pair; molecular orbital reasoning}]
Which of the following is incorrect
\begin{qoptions}
\opt{A}{$\mathrm{H_2>H}$ (ionization energy)}
\opt{B}{$\mathrm{O_2<O}$ (ionization energy)}
\opt{C}{$\mathrm{He<He^+}$ (ionization energy)}
\opt{D}{$\mathrm{N_2<N_2^+}$ (bond order)}
\end{qoptions}
\end{bankq}

\section*{Star Review (Other) papers}
\begin{bankq}[id=STAR01-P1-CHEM-S1-Q04, source={Star Review (Other) - Test 01 (12 Sep 2022)}, paper=1, qno=4, type=multi, subtopic={Periodic properties}, answer={C,D}, answer_src=key, difficulty=3, diff_parts={K1 S1 I1 C0 F1}, diff_note={concepts: electron affinity and ionisation energy trends | steps: compare each pair; exceptions in electron affinity}]
Which of the following set(s) have I$^{\mathrm{st}}$ element with both EA \& IE higher than 2$^{\mathrm{nd}}$ element.
\begin{qoptions}
\opt{A}{(N, O)}
\opt{B}{(P, Si)}
\opt{C}{(F, O)}
\opt{D}{(S, Se)}
\end{qoptions}
\end{bankq}
\begin{bankq}[id=STAR01-P2-CHEM-S1-Q09, source={Star Review (Other) - Test 01 (12 Sep 2022)}, paper=2, qno=9, type=single, subtopic={Trends}, answer={C}, answer_src=key, difficulty=1, diff_parts={K1 S1 I0 C0 F0}, diff_note={concepts: facts about Mendeleev table | steps: check each statement}]
Which of the following statement about Mendeleev's table is CORRECT?
\begin{qoptions}
\opt{A}{VIII$^{\mathrm{th}}$ group has 9 elements which are further divided in A \& B subgroups}
\opt{B}{Prediction of properties of Ge were done by leaving a gap below Mn}
\opt{C}{Cu and Cs were placed in same group}
\opt{D}{Correction in Atomic weight were done for B, In, U and Pt}
\end{qoptions}
\end{bankq}
\begin{bankq}[id=STAR01-P2-CHEM-S1-Q10, source={Star Review (Other) - Test 01 (12 Sep 2022)}, paper=2, qno=10, type=single, subtopic={Periodic properties}, answer={A}, answer_src=key, difficulty=3, diff_parts={K1 S2 I1 C0 F0}, diff_note={concepts: size of ions | steps: isoelectronic comparison; group trend; identify the wrong order}]
Which of the following ionic size order is incorrect.
\begin{qoptions}
\opt{A}{$\mathrm{Se^{2-}>I^-}$}
\opt{B}{$\mathrm{S^{2-}>N^{3-}}$}
\opt{C}{$\mathrm{C^{4-}>P^{3-}}$}
\opt{D}{$\mathrm{K^+>Sr^{2+}}$}
\end{qoptions}
\end{bankq}
\begin{bankq}[id=STAR01-P2-CHEM-S1-Q18, source={Star Review (Other) - Test 01 (12 Sep 2022)}, paper=2, qno=18, type=single, subtopic={Periodic properties}, answer={B}, answer_src=key, difficulty=2, diff_parts={K1 S1 I1 C0 F0}, diff_note={concepts: second ionisation energies | steps: configuration of the cations; order}]
Second ionsiation potential of Li, Be, B is in the order :
\begin{qoptions}
\opt{A}{Li $>$ Be $>$ B}
\opt{B}{Li $>$ B $>$ Be}
\opt{C}{Be $>$ Li $>$ B}
\opt{D}{Be $>$ B $>$ Li}
\end{qoptions}
\end{bankq}
\begin{bankq}[id=STAR01-P2-CHEM-S1-Q19, source={Star Review (Other) - Test 01 (12 Sep 2022)}, paper=2, qno=19, type=single, subtopic={Periodic properties}, answer={A}, answer_src=key, difficulty=3, diff_parts={K1 S2 I1 C0 F0}, diff_note={concepts: trends in ionisation energy and electron affinity | steps: compare nitrogen and oxygen; compare oxygen and sulphur; electronegativity relation}]
Consider the following statements.\par
(P) IE (I) of nitrogen atom is more than IE(I) of oxygen atom.\par
(Q) Electron affinity of oxygen is less than sulphur atom.\par
(R) Electronegativity of an atom has no relation with its ionization enthalpy and electron gain enthalpy\par
Which of the above statements are correct?
\begin{qoptions}
\opt{A}{P, Q}
\opt{B}{P, R}
\opt{C}{Q, R}
\opt{D}{P, Q and R}
\end{qoptions}
\end{bankq}
\begin{bankq}[id=STAR01-P2-CHEM-S2-Q08, source={Star Review (Other) - Test 01 (12 Sep 2022)}, paper=2, qno=8, type=numerical, subtopic={Periodic properties}, answer={3}, answer_src=solved, working={$X=3+6=9$, $Y^+$ has 18 electrons in the outer shell so $Y=\mathrm{Cu}$. True orders: Zn larger than Cu, IE1 of Cu above Ag, $\mathrm{Cu_2O}$ more basic than ZnO; the other two are false, so 3}, difficulty=7, diff_parts={K1 S3 I2 C1 F1}, diff_note={concepts: oxidation states and trends of an element | steps: oxidation number of iron; of chromium in the peroxide; identify Y; check each order; count}]
The sum of oxidation state of Fe in $\mathrm{K_3[Fe(CN)_6]}$ and Cr in $\mathrm{CrO_5}$ is "$X$". Monopositive ion of an element "$Y$" has exactly "$2X$" number of electrons in ultimate shell. If the element belongs to 4$^{\mathrm{th}}$ period then number of correct orders among the following is \_\_\_\_\_\_\_.\par
(i) Zn $>$ Y (Atomic size)\par
(ii) Y $>$ Zn ($\mathrm{IE_1}$)\par
(iii) Y $>$ Ag ($\mathrm{IE_1}$)\par
(iv) Y $>$ Na (number of unpaired electrons)\par
(v) $\mathrm{Y_2O>ZnO}$ (basic nature)
\end{bankq}
\begin{bankq}[id=STAR02-P2-CHEM-S3-Q06, source={Star Review (Other) - Test 02 (10 Oct 2022)}, paper=2, qno=6, type=numerical, subtopic={Trends}, answer={4}, answer_src=key, difficulty=7, diff_parts={K1 S3 I2 C1 F1}, diff_note={concepts: nine property orders | steps: each order in turn; acidity; boiling point; bond energy; proton affinity; density; count correct}]
The number of correct orders wrt the property mentioned along with is:\par
(i) $\mathrm{HF>HCl>HBr>HI}$ (Boiling point)\par
(ii) $\mathrm{H_2S<H_2Se<H_2Te}$ (Acidic strength)\par
(iii) $\mathrm{Cl_2>F_2}$ (Bond energy)\par
(iv) $\mathrm{NH_3>PH_3>AsH_3>SbH_3}$ (Proton affinity)\par
(v) $\mathrm{MnO_2>Mn_2O_7}$ (Acidic strength)\par
(vi) $\mathrm{HOCl>HClO_3}$ (Acidic strength)\par
(vii) $\mathrm{H_2O(s)<D_2O(s)}$ (Density)\par
(viii) $\mathrm{CH_3CH_2OH<CH_3CHO<CH_3COOH}$ (Boiling point)\par
(ix) $\mathrm{H_2O_2>H_2O}$ (Strength of H-bonding)
\end{bankq}

\section*{Nurture (General) papers}
\begin{bankq}[id=NUR04-P1-CHEM-S2-Q11, source={Nurture (General) - Test 04 (11 Feb 2024)}, paper=1, qno=11, type=multi, subtopic={Periodic properties}, answer={A,B,C}, answer_src=solved, working={O- + e, Ar ionisation and Ar + e are endothermic; H + e is exothermic}, difficulty=4, diff_parts={K1 S2 I1 C0 F1}, diff_note={concepts: endothermic electron gain or loss | steps: ionisation; electron gain of noble gas; second electron gain}]
Choose the process which is/are endothermic.
\begin{qoptions}
\opt{A}{O$^{-1}$ + e$^-\longrightarrow$ O$^{-2}$}
\opt{B}{Ar$\longrightarrow$ Ar$^+_{\,}$ + e$^-$}
\opt{C}{Ar + e$^-\longrightarrow$ Ar$^-$}
\opt{D}{H + e$^-\longrightarrow$ H$^-$}
\end{qoptions}
\end{bankq}
\end{document}
